\documentclass[11pt]{article}
\pdfinfoomitdate=1
\pdftrailerid{}
\pdfsuppressptexinfo=15
\usepackage[T1]{fontenc}
\usepackage{lmodern}
\usepackage[margin=1in]{geometry}
\usepackage{amsmath,amssymb,amsthm,mathtools,microtype}
\usepackage{enumitem,aliascnt,needspace}
\PassOptionsToPackage{hyphens}{url}
\usepackage[hidelinks]{hyperref}
\usepackage[nameinlink,noabbrev]{cleveref}
\numberwithin{equation}{section}
\newtheorem{theorem}{Theorem}[section]
\newaliascnt{proposition}{theorem}
\newtheorem{proposition}[proposition]{Proposition}
\aliascntresetthe{proposition}
\crefname{proposition}{proposition}{propositions}
\Crefname{proposition}{Proposition}{Propositions}
\newaliascnt{lemma}{theorem}
\newtheorem{lemma}[lemma]{Lemma}
\aliascntresetthe{lemma}
\crefname{lemma}{lemma}{lemmas}
\Crefname{lemma}{Lemma}{Lemmas}
\newaliascnt{corollary}{theorem}
\newtheorem{corollary}[corollary]{Corollary}
\aliascntresetthe{corollary}
\crefname{corollary}{corollary}{corollaries}
\Crefname{corollary}{Corollary}{Corollaries}
\newaliascnt{assumption}{theorem}
\newtheorem{assumption}[assumption]{Assumption}
\aliascntresetthe{assumption}
\crefname{assumption}{assumption}{assumptions}
\Crefname{assumption}{Assumption}{Assumptions}
\theoremstyle{remark}
\newaliascnt{remark}{theorem}
\newtheorem{remark}[remark]{Remark}
\aliascntresetthe{remark}
\crefname{remark}{remark}{remarks}
\Crefname{remark}{Remark}{Remarks}
\newcommand{\R}{\mathbb R}
\newcommand{\E}{\mathbb E}
\newcommand{\Pp}{\mathbb P}
\newcommand{\eps}{\varepsilon}
\newcommand{\dd}{\,\mathrm d}
\newcommand{\ind}{\mathbf1}
\newcommand{\TV}{\mathrm{TV}}
\newcommand{\Cov}{\operatorname{Cov}}
\newcommand{\Var}{\operatorname{Var}}
\newcommand{\osc}{\operatorname{osc}}
\newcommand{\kap}{\kappa}
\newcommand{\logp}{\log_+}
\newcommand{\sourcelabel}[1]{\nolinkurl{#1}}
\newif\iffunctional
\functionaltrue
\title{Functional hard-sphere fluctuations\\on regular kinetic intervals}
\author{}
\date{Research draft, version 1\\7 October 2026}
\hypersetup{pdftitle={Functional hard-sphere fluctuations on regular kinetic intervals},pdfsubject={Research draft, version 1 with explicit inherited analytic dependencies},pdfauthor={}}
\begin{document}
\maketitle
\begin{abstract}
We prove functional Gaussian convergence of exactly centered hard-sphere
fluctuations throughout a prescribed regular Boltzmann interval.
The limit is a Radon law on the path space of tempered distributions,
with the generalized Skorokhod topology induced by the strong dual,
and is supported on strongly continuous paths. We use the explicitly stated analytic
history package and finite-dimensional limit of the pinned upstream
manuscript. New quantitative first-recollision estimates control connected
cyclic histories by $C_h\mu^{1-h}\sqrt{\eps(1+\log(1/\eps))}$,
where $\mu=\eps^{-2}$. A joint squared-flux argument preserves the source's
factorial summability, including all permitted later weights.
Actual pasted-law increment cumulants follow after microscopic remainder
inversion. Passive collision records and a calendar-shift comparison give
a microscopic modulus, closing a two-grid tightness argument. The true flow
is reached by uniform exact-centered path comparison. No quantitative rate
of convergence of the fluctuation law or its limiting covariance is asserted.
\end{abstract}
\paragraph{Dependency and manuscript status.}
The analytic history construction, its coarse estimates, and the existing
finite-dimensional Gaussian limit are imported from the OpenAI
preprint~\cite{U}, at the exact snapshot specified below.
The new Gaussian geometry and weighted summation are proved here.
This is a research draft (version 1); external peer review and a comprehensive
publication-priority review have not been completed.
The strong-dual path topology is specified explicitly; convergence in a fixed negative-Sobolev path norm is not asserted. No publication-priority claim is made.
\section{Model and results}\label{sec:model}
Write $z=(x,v)\in\R^3\times\R^3$. Define
\[
 \|g\|_{\mathrm{Bol},\beta}
 =\sum_{a\in\mathbb Z^3}\sup_{|x-a|\leq1,\,v\in\R^3}
       e^{\beta|v|^2}|g(x,v)|.
\]
\begin{assumption}\label{ass:kinetic}
The initial datum $f_0$ is a smooth nonnegative probability density with
$\|f_0\|_{\mathrm{Bol},2\beta}
+\|\nabla_xf_0\|_{\mathrm{Bol},2\beta}<\infty$.
On a fixed interval $[0,T]$ there is a nonnegative classical Boltzmann
solution $f$ with $f|_{t=0}=f_0$ satisfying
$\sup_{t\leq T,x,v}e^{2\beta|v|^2}f_t(x,v)<\infty$.
\end{assumption}
The scattering convention is
\[
 B(v-v_*,\omega)=[(v-v_*)\cdot\omega]_+,\quad
 v'=v-[(v-v_*)\cdot\omega]\omega,\quad
 v_*'=v_*+[(v-v_*)\cdot\omega]\omega ,
\]
with ordinary surface measure on $S^2$, and
$Q(f,f)=\int B(f'f_*'-ff_*)\dd v_*\dd\omega$.
For diameter $\eps$ put $\mu=\eps^{-2}$ and use the grand-canonical law
\[
 \Pp_\eps(N,\dd Z_N)=
 Z_\eps^{-1}\frac{\mu^N}{N!}
 \prod_{j=1}^Nf_0(z_j)
 \prod_{i<j}\ind_{\{|x_i-x_j|>\eps\}}\dd Z_N .
\]
The true flow has elastic hard-sphere collisions. We also use the single
restarted pasted flow of~\cite[\S3]{U}: in each fixed slab, a proposed
collision is retained only while its connected component remains below
the prescribed size and cycle-excess caps. The parameter $\Gamma$ and
the number $L$ of slabs are fixed before choosing any moment order.
The size cutoffs are $A_L=|\log\eps|$, $\Lambda_i=A_i^{30}$,
and $A_{i-1}=\Lambda_i^{30}$, with harmless integer rounding.
A proposed C contact is accepted only if the new component has at most
$\Lambda_i$ labels and excess at most $\Gamma$, with repeated edges counted.
Otherwise the lines continue straight. Restarting changes only prehistory,
not positions or velocities; overlaps are allowed for this auxiliary flow.
Set $\tau=T/L$ and, for real $\phi\in C_c^\infty(\R^6)$,
\[
 \zeta_t^{\eps,p}(\phi)=\mu^{-1/2}
 \left(\sum_{i=1}^N\phi(z_i^p(t))
             -\E\sum_{i=1}^N\phi(z_i^p(t))\right).
\]
At every collision time we take the post-collision, right-continuous
velocity version, for both true and pasted flows. Positions are continuous.
The imported fixed-time laws are unchanged by this convention, as justified
below. Every centering is the exact expectation of the indicated microscopic law.

\begin{lemma}[Path convention]\label{lem:convention}
Replacing incoming values at collision instants by outgoing values gives
cadlag empirical test paths and changes neither a deterministic-time law
nor its expectation.
\end{lemma}
\begin{proof}
Conditional on each finite N, the initial law is absolutely continuous.
The fixed-time collision sets are Lebesgue-null under the hard-sphere and
capped flows; the latter preserve absolute continuity through their
countable injective volume-preserving branches
\cite[Appendix A, proof of Proposition 4.5, items (1) and (5)]{DHM}.
Induct over the fixed slab restarts for the pasted law. At each fixed time
both conventions consequently agree almost surely. A finite list of fixed
times has the same joint law. The almost-everywhere flows have the usual
one-sided collision limits; assigning the post-collision values gives
cadlag paths. For bounded tests the sum is dominated by a constant times N,
whose expectation is finite, so the mean is cadlag by dominated convergence.
\end{proof}
Put $r_\eps=\sqrt{\eps(1+\log(1/\eps))}$ for $0<\eps\leq e^{-1}$.

\begin{theorem}[Quantitative cyclic histories]\label{thm:cycles}
Under \Cref{ass:kinetic} and the imported package in
\Cref{prop:imports}, fix a factorial order $h$, a bounded finite record
list, and a complexity cutoff $K$.
The absolute sum of retained connected histories with trivial initial
cumulants and at least one redundant contact is at most
$C_{K,h,T}\mu^{1-h}r_\eps$.
The sum includes all sizes, all slab placements, and all permitted
post-first-return weights. The repeated-earlier-edge sector has the
stronger bound $C_{K,h,T}\mu^{1-h}\sqrt\eps$.
One sufficiently fine fixed slab decomposition works at every fixed
$K,h$.
\end{theorem}

\begin{theorem}[Microscopic increment cumulants]\label{thm:increments}
For the genuine pasted law and every fixed integer $k\geq2$,
\[
 \left|\kap_k\bigl(\zeta_t^{\eps,p}(\phi)
                      -\zeta_s^{\eps,p}(\phi)\bigr)\right|
 \leq C_{k,\phi}\eps^{k-2}(|t-s|+r_\eps),
 \qquad 0\leq s\leq t\leq T.
\]
The same dynamics is used at all orders. The sufficiently small upper
bound on $\eps$ may depend on the fixed order.
In particular $r_\eps\leq C_\eta\eps^\eta$ for every fixed $\eta<1/2$.
\end{theorem}

Bodineau, Gallagher, Saint-Raymond and Simonella~\cite{BGSS} already proved
nonequilibrium Gaussian-process convergence on a short time interval.
Their later work~\cite{BGSSeq} treats arbitrarily long times at equilibrium.
The source theorem~\cite{U} treats a regular kinetic lifespan away from
equilibrium, but explicitly concerns finite-dimensional laws and does not
assert path-space tightness. Gaussian-process convergence itself is
therefore not the novelty at issue. The functional extension below concerns
the source's precise regular-lifespan hypotheses.

\paragraph{Strong distribution-valued conclusion.}
The principal functional result, \Cref{thm:tempered}, gives weak convergence
of the centered empirical signed measures in
$D([0,T],\mathcal S'_\beta(\R^6))$, for the generalized $J_1$ topology
induced by the strong dual of the Schwartz space. Its limit is Radon and
strongly continuous. The proof includes the extension to every Schwartz
test and the path-law measurability and identification needed for the
nuclear-space tightness theorem. No rate of convergence of these path laws
or of their covariance to the Boltzmann limit is claimed.

\section{The imported history package}\label{sec:imports}
We state the analytical and combinatorial inputs from~\cite{U}, always
at the exact commit specified in the bibliography.
Their proofs, including the positive history estimates and kinetic
construction based on~\cite{DHM}, are not reproduced. Our results are
conditional on this inherited package; we add no equilibrium,
small-data, or shortened-lifespan assumption.

For a fixed bounded record attached to each particle, write $\xi_i$ for
its phase point together with that record. The ordered normalized
factorial measure of the actual pasted law is
\[
 F_h^\eps=\mu^{-h}\E\sum_{i_1,\ldots,i_h\ \text{distinct}}
           \delta_{(\xi_{i_1},\ldots,\xi_{i_h})}.
\]
Its factorial cumulants are the finite partition polynomials
\[
c_h^\eps=\sum_{\pi\in\mathcal P([h])}
(-1)^{|\pi|-1}(|\pi|-1)!\bigotimes_{B\in\pi}F_{|B|}^\eps,
\]
with tensor factors assigned to their labeled slots. A hat denotes the
source's finite retained hierarchy and its corresponding cumulants.
The source retains layered C/O histories with the cluster caps above,
forest O-graphs between C-components, rooted layer components, forest
initial links, and the continuation rules of its Definition
\sourcelabel{def:dhm-history}. Its complexity is
$\rho=|H_0|+\sum_i(|H_i|+r_i)$: $H_i$ is the set of designated or
continuing roots in slab $i$, $H_0$ consists of lines incident to initial
links, and $r_i$ is the slab's cycle excess. ``Fixed complexity'' means
$\rho\leq K$. Thus it does not bound the total particle count.
A global vertex is one physical particle line. A continued line retains
its old phase coordinate; a new interface birth has an independent one.
A contact between previously disjoint chronological components is a
merger; all others are redundant. C scatters and O continues straight.
At complexity $K$, let $m$ count birth lines, $n_i$ lines in slab $i$,
$a_i$ its mergers, $q$ interface continuations and $p$ physical initial
links. For a connected history,
\begin{equation}\label{eq:counts}
 a=\sum_i a_i=m-1-p,\qquad n=\sum_i n_i=m+q,\qquad
 \sum_i(n_i-a_i)=q+1+p=O_K(1).
\end{equation}
The redundant-contact count is also $O_K(1)$.

\Needspace{6\baselineskip}
\begin{proposition}[Imported package]\label{prop:imports}
The following hold for the single pasted dynamics.
\begin{enumerate}[label=(\roman*)]
\item Independent birth inputs have one common Gaussian Bol envelope.
Integrating a common spatial translation per component and taking spatial
suprema of the other inputs gives $C^m e^{-bE}$, uniformly in every
displacement, where $E$ is total birth energy. Active energy never exceeds
$E$.
\item Merger coordinates consist of birth velocities, representative
translations, an initial displacement per physical link, and a free
$(t,\omega)$ per dynamic merger. Its Jacobian is $\eps^2B\dd t\dd\omega$.
A redundant incoming intersection has indicator mass and a determined
velocity update. A missing branch is extended by zero. Every step uses
only preceding parameters.
\item The normalization is
$\mu^{m-h}\eps^{2a}\eps^{3p}=\mu^{1-h}\eps^p$.
Beyond ordered merger pair lists, choices cost
$C_K(1+n)^{d_K}C^a$. Fresh-label denominators give $\prod_i n_i!$ in the
denominator up to fixed polynomial factors. Merger pairs are not included
in that prefactor.
\item The actual marked factorial hierarchy and its retained version obey
\[
 \|F_h^\eps-\widehat F_h^\eps\|_{\TV}
       \leq C_h e^{-c_h\Lambda_L},\qquad
 \Lambda_L=(\log(1/\eps))^{30}.
\]
This is uniform at every fixed order and bounded record count, including
partial final slabs. It is unweighted total variation.
Retained factorial-cumulant inversion keeps exactly connected diagrams.
\item For every fixed order and prescribed $P$, a fixed complexity cutoff
makes the unscaled high-complexity retained tail $O(\eps^P)$.
Physical links gain $\eps^p$. Nonphysical initial tails can be erased and
assigned arbitrary powers by an auxiliary depth choice.
These choices do not change the microscopic dynamics.
\item Actual normalized factorial masses are at most one, and every
fixed moment of $N/\mu$ is uniformly bounded.
\item For the same dynamics, the actual unmarked cumulants obey
$\|c_{h,t}^\eps\|_{\TV}\leq C_h\mu^{1-h}$ at every fixed order.
The exactly centered pasted observations at any fixed finite list of
compact smooth tests and times converge to the Gaussian law described
in \Cref{sec:functional}. Its target is a continuous linear map
$C_0(\R^6)\to L^2$ with
$\sup_{t\leq T}\|G_t(\phi)\|_{L^2}\leq C\|\phi\|_\infty$.
On the common-initial-data coupling, the probability that the true and
pasted trajectories differ anywhere in $[0,T]$ is at most
$C\eps^{21/20}$.
\end{enumerate}
Exponential bases and the choice of $L$ are independent of subsequent
fixed orders and cutoffs; front constants and polynomial degrees may
depend on them.
\end{proposition}

The following exact source locators distinguish the inputs. Labels refer
to the TeX sources at the pinned snapshot, not to claims proved here.
\begin{itemize}[leftmargin=1.5em]\raggedright
\item Items (i)--(iii): \cite[\S3]{U},
\sourcelabel{thm:dhm-package}; \cite[\S4]{U},
\sourcelabel{prop:one-root}; and \cite[\S6]{U},
\sourcelabel{eq:sharp-input-bounds},
\sourcelabel{lem:sharp-merger-coordinates},
\sourcelabel{lem:sharp-factorials}, and
\sourcelabel{lem:sharp-summability}.
\item Items (iv)--(v): \cite[\S4]{U},
\sourcelabel{lem:autonomous-partition},
\sourcelabel{lem:one-slab-expansion},
\sourcelabel{lem:marked-regrouping},
\sourcelabel{lem:diagram-labels},
\sourcelabel{lem:stopped-variation}, and
\sourcelabel{thm:marked-expansion}; physical-link scaling and initial-tail
deletion are in \cite[\S6]{U},
\sourcelabel{eq:sharp-initial-link} and
\sourcelabel{eq:sharp-scaling}.
\item Item (vi): \cite[\S2]{U}, \sourcelabel{prop:initial}, especially
\sourcelabel{eq:initial-factorial-upper} and
\sourcelabel{eq:initial-N-moments}.
\item Item (vii): \cite[\S6]{U},
\sourcelabel{thm:sharp-cumulants}; \cite[\S7]{U},
\sourcelabel{prop:cumulant-clt}; \cite[\S2]{U},
\sourcelabel{prop:gaussian-target} and
\sourcelabel{eq:target-field-bound}; and \cite[\S5]{U},
\sourcelabel{prop:cutoff}. The exact-centering argument from
\cite[\S8]{U}, \sourcelabel{lem:center-transfer}, is reproduced below.
\end{itemize}
Required record lines continue to time zero.
The finite identities depend on bounded output kernels and record count,
not on a positive separation between observation times.
Our Gaussian reference has covariance $\sigma^2I$ per birth velocity,
obtained after a fixed reservation of the exponent in (i).
It is not a conditional law of the physical gas.

\section{Gaussian geometry before the first redundancy}\label{sec:geometry}

\begin{lemma}[Transverse Gaussian bound]\label{lem:smallball}
Let $(D,U)$ be jointly Gaussian in $\R^3\times\R^3$, with $U$ centered
and isotropic. If $\lambda=\lambda_{\max}\Cov(D\mid U)>0$, then
\[
 \Pp\{\operatorname{dist}(D,\R U)\leq\eps\}
       \leq C\min(1,\eps/\sqrt\lambda).
\]
\end{lemma}
\begin{proof}
Choose a top eigenvector $e$. Conditional on $U=u$, project onto
$w=P_{u^\perp}e/\sin\theta$, where $\theta$ is the angle of $u,e$.
The variance of $w\cdot D$ is at least $\lambda\sin^2\theta$.
A one-dimensional Gaussian density gives
$C\min(1,\eps/(\sqrt\lambda\sin\theta))$, independently of its mean.
The direction of $U$ is uniform and $\E(\sin\theta)^{-1}=\pi/2$.
The zero and parallel cases are null.
\end{proof}

\subsection{Repeated earlier pairs}
\begin{lemma}[An erased normal coordinate]\label{lem:repeated}
Fix a chronological merger forest, its normals and C/O flags.
Suppose A and B have a prescribed contact at $t_0$, and the prospective
first redundant contact repeats this pair. Let $t_1>t_0$ be the first
later C contact touching either endpoint, and $\ell=t_1-t_0$.
At the beginning of the tested subsequent free segment, write $D$ for
their relative displacement and $Y$ for both endpoint velocities. Then
\[
 \Cov(D\mid Y)\succeq\sigma^2\ell^2 n_1n_1^T,
 \qquad
 \Pp_\gamma\{\text{return}\}\leq C_\sigma\min(1,\eps/\ell).
\]
If no such C contact exists, a second incoming crossing is impossible.
\end{lemma}
\begin{proof}
A C collision swaps the two normal components and retains the transverse
ones. It is orthogonal in all birth Gaussian coordinates; the current
velocity array is therefore a product Gaussian at fixed free normals.
After A and B first belong to one component, every merger touching an
endpoint imports a disjoint component. Its incident normal scalar is
Gaussian and independent of the entire previous tagged history.
That imported component cannot feed a scalar back through another merger:
such a contact would already be redundant.

Suppose A has the first C update. With $A_0,B_0$ the velocities just after
the earlier pair contact, set $\xi=A_0\cdot n_1$. It is independent of
$P_{n_1^\perp}A_0,B_0$ and all later imported scalars. At $t_1$ it is
erased from A and never re-enters a tagged endpoint before the first
redundancy. Consequently
\[
 D=\ell n_1\xi+R,\qquad
 \xi\ \text{is independent of }(R,Y).
\]
The initial contact displacement is a deterministic $\eps$-normal
offset. This proves the covariance inequality. Updating B changes only
the sign. The final relative velocity is isotropic, and the proof of
\Cref{lem:smallball} gives the stated probability bound.
Without a C update the relative line remains straight: a C contact
separates, while an O contact has only the outgoing exit. Neither permits
a second incoming crossing.
\end{proof}

The conditioning is only on the two endpoint velocities, not on all
partner velocities, realizability, or collision fluxes. The erased scalar
was exported to a partner and can be revealed by those stronger
conditionings.

\subsection{Nonneighbor pairs}
\begin{lemma}[Rank one across a forest cut]\label{lem:rankone}
The $3$ by $3$ velocity-transfer block between distinct labels in a
prescribed forest of C/O contacts has rank at most one.
\end{lemma}
\begin{proof}
Cut an edge on their unique path. It is the only contact crossing that
cut. A C contact transmits one normal scalar across it; O transmits none.
Every later update is internal to one side, so dependence on the original
velocity continues to factor through that scalar. This does not require
the chronological order to be monotone along the spatial path.
\end{proof}

Fix nonneighbors A,B. Orient tree edges $e=(p_e,q_e)$, put
$b_e=e_{p_e}-e_{q_e}$, and write
$b=e_A-e_B=\sum_es_eb_e$ for their signed spatial path.
The actual contact equations give, at a time $t_*$ after the prefix,
\begin{equation}\label{eq:position}
 D_\eps=\eps\sum_es_e\omega_e+
   \sum_es_e\int_{t_e}^{t_*}(v_{p_e}(r)-v_{q_e}(r))\dd r .
\end{equation}
This telescopes positions; it makes no constrained initial positions
independent. At fixed normals all velocities are linear in the birth array
$V$. Write $D_\eps=\eps\xi+L(t)V$, $U=QV$, $QQ^T=2I$. Then
\begin{equation}\label{eq:regression}
 \Cov(D_\eps\mid U)=\sigma^2H(t)H(t)^T,\qquad
 H(t)=L(t)\left(I-\tfrac12Q^TQ\right).
\end{equation}
Velocity maps depend on the prescribed order and normals, not on the
contact times inside their ordered domains.

\begin{lemma}[A nonzero contact-time slope]\label{lem:slope}
Let $e_*$ be the earliest edge on the A--B spatial path.
The derivative of the residual coefficient $H$ with respect to its
contact time has operator norm at least one.
\end{lemma}
\begin{proof}
Before $e_*$ the sum of the already-contacted path incidences is zero;
just afterward it is $e_p-e_q$. Differentiating \eqref{eq:position} gives
the negative relative velocity immediately after $e_*$.
Remove this edge from the final tree. The two sides have independent
current Gaussian coordinates, and no later contact crosses the cut.
If $K_A,K_B$ transfer the current $p,q$ velocities to the final A,B
velocities, conditioning on both final velocities gives covariance
\[
 \sigma^2(2I-K_A^TK_A-K_B^TK_B).
\]
At least one transfer is off-diagonal, because the path has at least two
edges. It has rank at most one by \Cref{lem:rankone} and norm at most one
by orthogonality. On its kernel the displayed matrix has a direction of
variance at least $\sigma^2$. Conditioning only on the final difference
increases covariance. Equation \eqref{eq:regression} proves the claim.
\end{proof}

\begin{proposition}[Time-averaged nonneighbor event]\label{prop:nonneighbor}
Let $a_i$ count all free merger times in the full schedule, including
post-return mergers, and let $i_*$ contain $e_*$. The Gaussian prefix
return event $G$ satisfies
\[
 \int_{\prod_i\mathcal S_{a_i}(\tau)}\Pp_\gamma(G)\dd t
 \leq C\prod_i\frac{\tau^{a_i}}{a_i!}
 \min\left(1,\frac{a_{i_*}\eps}{\sigma\tau}
                  [1+\logp(\sigma\tau/\eps)]\right).
\]
This is uniform in all free normals and prescribed forest shapes.
\end{proposition}
\begin{proof}
Vary $t_{e_*}$ between its nearest full-schedule neighbors in its slab,
using slab endpoints if needed. Hold all other parameters fixed.
The slope lemma supplies unit vectors $w,z$ with
$w^TH(t)z=\alpha t_{e_*}+\beta$, $|\alpha|\geq1$.
The largest conditional variance is at least
$\sigma^2|\alpha t_{e_*}+\beta|^2$. For an interval $J$ of length $s$,
\[
 \int_J\min(1,\rho/|\alpha t+\beta|)\dd t
       \leq2\rho[1+\logp(s/(2\rho))],\qquad \rho=\eps/\sigma.
\]
The even decreasing integrand has maximal interval integral when centered
at its singularity. Apply \Cref{lem:smallball}, use $s\leq\tau$, and
remove the selected time from its ordered simplex. The remaining volume
is $\tau^{a_{i_*}-1}/(a_{i_*}-1)!$; every other slab factorial is unchanged.
Also use the trivial probability bound.

For fixed birth and ending interfaces, allocate all birth Gaussian
coordinates in advance, leave unborn coordinates unused, and freely
extend terminated lines only for \eqref{eq:position}. The auxiliary
extensions introduce no contact or random coordinate and cancel in the
spatial telescoping sum. Varying the selected contact within its own slab
retains the original interface assignments.
\end{proof}

An untruncated inverse covariance moment at exponent one is not claimed:
two O contacts on the path $1$--$2$--$3$ give covariance proportional to
the square of their time gap, whose inverse square root is not integrable.
The truncated probability estimate legitimately retains one logarithm.

\section{Fluxes and the complete retained size sum}\label{sec:summation}
Assume first that initial cumulants are trivial, so $p=0$.
Fix a bounded-complexity exceptional skeleton, including the first
redundant-contact slot. Let $s$ range over decorated merger pair schedules.
After spatial-envelope domination write $F_s$ for the full nonnegative
merger-flux weight, with all needed redundant-branch indicators retained.
It is zero on an unrealizable branch. Let $G_s$ be the Gaussian transverse
return demand determined solely by the formal forest before its first
redundancy. Then $\operatorname{supp}F_s\subset G_s$.
No post-return Gaussian invariance is asserted.

\begin{lemma}[Squared-flux sum]\label{lem:squared}
At fixed merger times,
\[
 \sum_s\int F_s^2\dd\Omega
       \leq C^aE^a\prod_i n_i^{a_i}.
\]
The number of ordered pair schedules is at most
$C^a\prod_i n_i^{2a_i}$.
\end{lemma}
\begin{proof}
At a valid branch node in slab $i$,
\[
 \sum_{p<q}|v_p-v_q|^2
 =n_i\sum_p|v_p|^2-\left|\sum_pv_p\right|^2\leq n_iE.
\]
The squared incoming flux is at most its squared relative speed.
Integrate the free normal and sum C/O choices, then induct backward
through the pair-choice tree. Every valid child preserves the energy
bound, including all intervening determined redundant scatterings.
Invalid children remain zero; adding other pair speeds only enlarges
the estimate. Thus no fictitious velocity is needed on an invalid branch.
The schedule count follows by permitting every unordered pair at every
merger, ignoring forest restrictions only in this numerical upper bound.
\end{proof}

Apply Cauchy--Schwarz jointly in schedules and the product parameter
measure $\nu$:
\begin{equation}\label{eq:jointCS}
 \sum_s\int F_s\dd\nu
 \leq\left(\sum_s\int F_s^2\dd\nu\right)^{1/2}
       \left(\sum_s\int\ind_{G_s}\dd\nu\right)^{1/2}.
\end{equation}
For $m$ birth Gaussians and $a=m-1\leq n$,
\[
 \E_\gamma E^a=(2\sigma^2)^a
       \frac{\Gamma(3m/2+a)}{\Gamma(3m/2)}
       \leq(C_bn)^a .
\]
In the nonneighbor case include \emph{all} free merger times in $\nu$.
With $V_t=\prod_i\tau^{a_i}/a_i!$ and
$B_\eps=1+\logp(\sigma\tau/\eps)$, the two factors in \eqref{eq:jointCS},
before taking square roots, are bounded by
\[
 C^{m+a}V_t n^a\prod_i n_i^{a_i},
 \qquad
 C^aV_t\prod_i n_i^{2a_i}\frac{Cn\eps B_\eps}{\sigma\tau}.
\]
Their product restores the full time volume and gives
\begin{equation}\label{eq:preStirling}
 C^{n+a}V_t n^{a/2}\prod_i n_i^{3a_i/2}
              \sqrt{\frac{n\eps B_\eps}{\sigma\tau}} .
\end{equation}
Only the second factor uses the Gaussian forest geometry. In particular
later merger weights and nonlinear determined returns stay in the first.

In the repeated-pair sector, refine the exceptional skeleton by the
earlier pair-contact slot and the first later C endpoint-update slot.
These at most $a^2\leq n^2$ choices are absorbed in the fixed-complexity
polynomial. Within each refined class $t_0$, $t_1$, and the elementary
gap selected below are common to all schedules. Classes with no such
C update vanish by \Cref{lem:repeated}.
We may now apply \eqref{eq:jointCS} at fixed times
and use \Cref{lem:repeated}. The resulting factor is
$\min(1,\sqrt{\eps/(\sigma\ell)})$.
If the earlier edge lies in slab $i$, choose either its next same-slab
merger gap or the terminal gap to that slab's endpoint. In both cases
the chosen gap $\Delta$ is at most $\ell$. Under the normalized ordered
slab simplex,
\[
 \E\Delta^{-1/2}
 =\tau^{-1/2}
   \frac{\Gamma(1/2)\Gamma(a_i+1)}{\Gamma(a_i+1/2)}
 \leq C\sqrt{(1+a_i)/\tau}.
\]
Thus \eqref{eq:preStirling} holds without $B_\eps$ in that sector.

\begin{proof}[Proof of \Cref{thm:cycles}]
Insert the fresh-label denominator $\prod_i n_i!$ and the fixed
complexity polynomial. The source's elementary Stirling calculation,
with $\delta_i=n_i-a_i\geq0$, gives
\[
 \frac{n_i^{3a_i/2}}{n_i!a_i!}
 \leq e^{2n_i}(1+n_i)^{\delta_i}n_i^{-n_i/2},
 \qquad
 \prod_i n_i^{n_i}\geq(n/L)^n .
\]
Because $\sum_i\delta_i=O_K(1)$, \eqref{eq:preStirling} is bounded by
\[
 C_{K,h,T}\mu^{1-h}r_\eps
       (1+n)^{d_K}(C_b\tau\sqrt L)^n .
\]
The activity factor is exactly $\mu^{m-h}\eps^{2a}=\mu^{1-h}$.
The number of slab-size vectors of total $n$ is polynomial for fixed $L$.
Choose $L$ once so that $C_bT/\sqrt L<1$. The size series converges.
The repeated sector gives $\sqrt\eps$ in place of $r_\eps$.
The first redundant pair is either an earlier edge or a nonneighbor,
so these cases exhaust the $p=0$ cyclic class.
\end{proof}

Physical initial-link sectors separately have $\eps^p$, $p\geq1$,
and the original summable majorant bounds them by
$C_K\mu^{1-h}\eps$. Nonphysical initial tails and high complexity
are treated by their actual arbitrary-power estimates in
\Cref{prop:imports}; they are not made independent by the new geometry.

\begin{remark}
Uniform $C^n$ flux moments for each fixed particle tree would be false.
For collinear C contacts $(1,2),(2,3),\ldots,(n-1,n)$, the same ancestral
normal scalar is transported along the chain. Its flux product is
$\prod_{j=2}^n(X_1-X_j)_+$ with independent Gaussian $X_j$.
The event $X_1/\sigma\in[2\sqrt{n-1},2\sqrt{n-1}+1]$ and
$|X_j|/\sigma\leq1$ gives expectation at least
$c^n\sigma^{n-1}(n-1)^{(n-1)/2}$.
The particle tree has degree two. Our all-schedule squared-flux proof
preserves, and then cancels, this Gaussian growth through the two
factorial denominators.
\end{remark}

\section{Actual microscopic increment cumulants}\label{sec:increments}
Fix $I=(s,t]$, and mark each designated particle by
$y=\phi(z_t)-\phi(z_s)$ using two bounded observation pulses.
Put $M=\|\phi\|_\infty$ and $D=\|v\cdot\nabla_x\phi\|_\infty$.
The record satisfies $|y|\leq2M$, and along every prescribed particle line
\begin{equation}\label{eq:incrementpath}
 |y|\leq D|I|+2M N_I,
\end{equation}
where $N_I$ may count all contacts of its diagram. Counting O contacts
only increases this bound.

\begin{lemma}[One-window tree localization]\label{lem:treewindow}
For connected retained merger trees of factorial order $h$, fixed
complexity and fixed total mark degree $k$, testing at least one increment
factor gives an absolute bound $C_k\mu^{1-h}|I|$.
\end{lemma}
\begin{proof}
For $a$ ordered times in a slab of length $\tau$,
\[
 \int_{\mathcal S_a(\tau)}
      \sum_{j=1}^a\ind_{\{t_j\in I\}}\dd t
       =\frac{a|I\cap\text{slab}|\tau^{a-1}}{a!}.
\]
This follows by symmetrizing over the cube. Before time integration the
source's pair-flux and Gaussian majorants are uniform in the merger
times. Inserting the count therefore costs only a polynomial size factor
times $|I|/\tau$, preserving summability.
Apply \eqref{eq:incrementpath} to one marked factor and bound the others
by $2M$. A contact-free term contributes only its streaming factor.
Required record lines continue to time zero, so no part of this trajectory
is hidden in an unmarked independent input.
\end{proof}

\begin{lemma}[Microscopic remainder after cumulant inversion]\label{lem:remainder}
For every fixed $k$ and all $h\leq k$,
\[
 \|c_h^\eps-\widehat c_h^\eps\|_{\TV}
       \leq C_k e^{-c_k\Lambda_L},
\]
uniformly in the two pulse times.
\end{lemma}
\begin{proof}
The actual normalized factorial masses are at most one.
Thus the retained factorial variations are at most one plus their
exponentially small errors. Invert the finite partition identity:
\[
 c_h=\sum_{\pi\in\mathcal P([h])}
       (-1)^{|\pi|-1}(|\pi|-1)!
       \bigotimes_{B\in\pi}F_{|B|}.
\]
Expand each tensor-product difference one factor at a time and use the
factorial mass bounds. This proves the claim.
Record updates are measurable pushforwards of norm at most one; their
number and bounded ranges, not their time separation, control the source
estimates. No velocity weight is attached to this error.
\end{proof}

\begin{proof}[Proof of \Cref{thm:increments}]
Fix $k$ and choose one complexity cutoff covering all $h\leq k$ with
unscaled tail $O(\eps^{2k+2})$. Choose the initial-tail depth afterward.
Neither operation changes the dynamics.
For integers $b_j\geq1$ with $\sum_jb_j=k$, split retained connected
diagrams into merger trees, cyclic $p=0$ histories, physical-link
histories and the initial-tail remainder.
\Cref{lem:treewindow,thm:cycles} and the separate link gain imply
\[
 \left|\int c_h^\eps\prod_{j=1}^h y_j^{b_j}\right|
       \leq C_k\mu^{1-h}(|I|+r_\eps).
\]
Here the actual microscopic error from \Cref{lem:remainder} is smaller
than every fixed power, and $\eps^{2k+2}$ is smaller than
$\mu^{1-h}r_\eps$ for $h\leq k$.

For $S_I=\sum_i y_i$, the finite marked partition identity is
\[
 \kap_k(S_I)=\sum_{\pi\in\mathcal P([k])}
      \mu^{|\pi|}\int c_{|\pi|}^\eps
                         \prod_{B\in\pi}y_B^{|B|}.
\]
It follows by expanding
$e^{\lambda S_I}=\prod_i[1+(e^{\lambda y_i}-1)]$ and taking the logarithm;
through degree $k$ it is a finite formal identity.
At fixed $\eps$ it is also justified by the exponential moments of $N$.
For $k\geq2$, exact centering changes no cumulant and each partition
has normalization
$\mu^{-k/2}\mu^{|\pi|}\mu^{1-|\pi|}
=\mu^{1-k/2}=\eps^{k-2}$.
This proves the theorem, including uniformity over pulse times.
The size-summation bases are independent of fixed order, so $\Gamma,L$
remain the same throughout.
\end{proof}

\begin{corollary}[Even moments]\label{cor:moments}
For each fixed integer $p\geq1$, with $B_I=|I|+r_\eps$,
\[
 \E|\Delta_I\zeta^{\eps,p}(\phi)|^{2p}
       \leq C_p(B_I^p+\eps^{2p-2}B_I).
\]
If $0<a<1/2$ and $|I|\leq\eps^a$, the bound is $C_p\eps^{ap}$.
\end{corollary}
\begin{proof}
A centered moment partition has no singleton blocks.
A partition into $r\leq p$ blocks contributes
$C_p\eps^{2p-2r}B_I^r$. This geometric progression is bounded by a
constant times its two endpoint terms. For the last assertion,
$r_\eps=o(\eps^a)$ and
$2p-2+a\geq ap$.
\end{proof}

This argument proves a result for the actual pasted probability law,
not only a signed retained expansion. It does not transfer its high
moments to the true flow. The latter transfer will use uniform path
comparison, which requires substantially less.

\iffunctional
\section{Collision counts and microscopic windows}\label{sec:micro}
We prove the needed source extensions explicitly. In this section a
constant may depend on a requested fixed remainder accuracy.

\begin{lemma}[A passive incidence record]\label{lem:record}
The exact component expansion and unweighted variation estimates in
\Cref{prop:imports} also hold when a designated particle records its
number of accepted C contacts in an interval $I$.
\end{lemma}
\begin{proof}
In slab $i$ an accepted component has at most
$\Lambda_i-1+\Gamma$ edges. Thus each particle's total incidence record
is at most $B_\eps=\sum_i(\Lambda_i-1+\Gamma)$, a fixed power of
$|\log\eps|$. Normalize it by $B_\eps$.
In the source's autonomous partition identity, compatible autonomous
components reproduce the full accepted root history, hence this statistic.
The one-slab root-output kernel can therefore append it as one passive
coordinate. Spectator cancellation and the fresh-label identity are
unchanged. Pushforward has variation norm at most one; the crude operator
bound and stopped forest domination do not require derivatives or a
particular choice of output statistic.
Required record labels still continue to zero, while unmarked one-root
blocks are the same inputs $g_i$. Thus the source's finite regrouping and
coarse estimates apply with one bounded coordinate per required label,
not one coordinate per mesh point.
At total degree $k$, restoring the raw record costs $B_\eps^k$ on errors.
This is polylogarithmic, so exponentially small errors remain smaller than
every power, and arbitrary-power tails can be chosen after $k$.
On a retained diagram, raw mark powers cost only a fixed power of its
contact count, hence a fixed polynomial in size at complexity $K$.
\end{proof}

\begin{lemma}[Calendar localization]\label{lem:calendar}
For fixed factorial order, complexity and total contact-mark degree, the
sum of connected retained absolute marked integrals is at most
$C\mu^{1-h}|I|$, provided at least one marked contact is required to lie
in $I$. Nonphysical initial tails may instead be assigned arbitrary-power
errors.
\end{lemma}
\begin{proof}
First integrate out every common spatial translation and replace the
birth factors by their uniform envelope $C^m e^{-bE}$.
The remaining quotient coordinates are $(V,\Omega,t,y)$; there is no
remaining integration in an unrestricted translation variable.
Collapse identity interface gluings of continued lines.
For the positive comparison family, freely extend each global line
before its first and after its last prescribed contact. Redundant-contact
segments end at neighboring prescribed contacts or at infinity, not at
fixed physical lifetime interfaces. Keep the full chronological slot
order and the requirement that each redundant incoming intersection
exists; otherwise its branch weight is zero.

For $u\geq0$ set
\[
 t_e\longmapsto t_e+u,\qquad
 y_{pq}\longmapsto y_{pq}-u(v_p^{\rm birth}-v_q^{\rm birth}),
\]
leaving birth velocities and free normals fixed.
This quotient-coordinate map has determinant one. Reconstructing virtual
paths inductively shows that the new path at time $t$ is the old path at
$t-u$, up to the already integrated representative translation.
Merger normals, velocities and fluxes are unchanged. A redundant
incoming intersection shifts by $u$ with its normal and update unchanged.
The first shifted contact is later than the original one, so its fixed
birth anchor remains on its incoming free segment.
Endings and nominal slab interfaces are only constraints in the original
domain; virtual comparison paths may extend beyond them.
All global coordinates remain in the energy accounting, and every update
preserves total birth energy.

Put all absolute birth, ending, interface and nominal-slab restrictions,
including those on determined redundant times, into a domain $D$.
The Gaussian, fluxes and virtual chronological branch indicators form
a calendar-invariant core $W$.
If $t_*$ is a selected contact, $t_*(\Phi_u\theta)=t_*(\theta)+u$.
Changing variables, enlarging all shifted domains to $D_\ell^+$, and
averaging over $0\leq u\leq\ell$ gives
\begin{equation}\label{eq:calendar-average}
 \int_D\ind_{\{t_*\in I\}}W\dd\theta
 \leq\frac{|I|}{\ell}\int_{D_\ell^+}W^+\dd\theta ,
\end{equation}
because $\int_0^\ell\ind_{\{t_*\in I+u\}}\dd u\leq|I|$.
The changed interval $I+u$ is essential.

If $p=0$, take $\ell=\tau$. Nominal merger-time intervals have length at
most $2\tau$, and remain ordered within each slab.
Their volumes are $(2\tau)^{a_i}/a_i!$.
Retain the original global slot order until the final volume upper bound.
The same pair-speed, Gaussian-energy and factorial estimates apply to
these virtual histories: the subset of $n_i$ nominal labels has
$\sum_{p<q}|v_p-v_q|\leq n_i^{3/2}\sqrt E$.
Thus their size majorant is
$C_K\mu^{1-h}(1+n)^{d_K}(2C_b\tau\sqrt L)^n$.
Taking $L$ sufficiently large once absorbs the enlargement.

If $p\geq1$, take $\ell=\eps<\tau$. An initial-link ball enlarges to
radius $\eps(R_{K,D}+\sqrt{2E})$. The $p$ link volumes are at most
$C_{K,D}\eps^{3p}(1+E)^{3p/2}$.
The normalization remains $\mu^{1-h}\eps^p$.
Averaging in \eqref{eq:calendar-average} pays $\eps^{-1}$; the link gain
leaves $\eps^{p-1}\leq1$.
Absorb the fixed energy polynomial into half the birth Gaussian:
$(1+E)^{3p/2}e^{-bE}\leq C_{b,p}e^{-bE/2}$.
The fraction reserved is fixed independently of $p$, so the exponential
size base stays independent of complexity and moment order.
Nonphysical initial tails are erased with their inherited power cost
and chosen afterward.
Selecting a contact and restoring raw mark powers add only polynomial
size factors. Summing the geometric size tail proves the lemma.
\end{proof}

\begin{proposition}[Actual incidence cumulants]\label{prop:counts}
Let $A_I$ be the total number of collision incidences, two per accepted
C contact. For every fixed integers $k\geq1$ and fixed accuracy $M>0$,
\[
 |\kap_k(A_I)|\leq C_{k,M}(\mu|I|+\eps^M)
\]
uniformly in $I\subset[0,T]$.
\end{proposition}
\begin{proof}
Apply \Cref{lem:record}. In a connected retained term of total mark degree
$k$, a nonzero mark product selects at least one C contact in $I$.
\Cref{lem:calendar} bounds its factorial-order-$h$ contribution by
$C_{h,K,k}\mu^{1-h}|I|$. The ordinary-cumulant partition identity then
multiplies by $\mu^h$, giving $C\mu|I|$.
Choose complexity, initial tails and stopped-error powers after $k,M$,
allowing for the fixed factors $\mu^h$ and $B_\eps^k$.
The actual hierarchy-to-cumulant remainder is handled exactly as in
\Cref{lem:remainder}. The leading constant may depend on $M$ through $K$;
there is no assertion of a constant uniform over all accuracy choices.
\end{proof}

\begin{proposition}[Microscopic modulus]\label{prop:micro}
For every fixed $J\geq2$, $\lambda>0$ and compact smooth $\phi$,
\[
 \Pp\left\{\sup_{|t-s|\leq\eps^J}
       |\zeta_t^{\eps,p}(\phi)-\zeta_s^{\eps,p}(\phi)|>\lambda\right\}
       \leq C_{\phi,\lambda,J}\eps^2
\]
for sufficiently small $\eps$.
\end{proposition}
\begin{proof}
Put $\delta=\eps^J$ and cover the interval by $O(T/\delta)$ overlapping
windows of length $2\delta$. The raw fourth-moment identity and
\Cref{prop:counts}, with the accuracy fixed sufficiently large after $J$,
give
\[
 \E A_I^4\leq C\sum_{r=1}^4(\mu|I|)^r+O(\eps^M).
\]
For example accuracy $M+6$ in each cumulant absorbs the other at most
three factors of order $\mu$. Markov and the grid union bound yield
\[
 \Pp\{\max_I\eps A_I>\lambda\}
 \leq C_\lambda(\eps^2+\delta+\eps^{-2}\delta^2+
                         \eps^{-4}\delta^3)
       +O(\eps^{M+4-J})=O(\eps^2).
\]
Pathwise, with $D_\phi=\|v\cdot\nabla_x\phi\|_\infty$,
\[
 \osc_I\left(\eps\sum_i\phi(z_i^p(t))\right)
       \leq\eps D_\phi N|I|+2\|\phi\|_\infty\eps A_I.
\]
There are no state jumps when the pasted prehistory restarts.
The single global $N$ term needs no grid union:
$\E N^2\leq C\mu^2$ bounds its probability by
$C\eps^2\mu^2\delta^2=C\eps^{2J-2}\leq C\eps^2$.
The exact mean increment is deterministically at most
$C_\phi\eps^{-1}\delta+O(\eps^M)$ by the first incidence cumulant.
It is smaller than a fixed threshold for small $\eps$.
This proves the statement, with constants allowed to depend on $J$.
\end{proof}

\section{Finite families of fluctuation processes}\label{sec:functional}
For a test $\phi$ write
$\Delta\phi=\phi(x,v')+\phi(x,v_*')-\phi(x,v)-\phi(x,v_*)$.
The inherited Gaussian finite-dimensional law has weak drift and noise
\begin{align*}
 A_t\phi&=v\cdot\nabla_x\phi+
       \int B(v-v_*,\omega)f_t(x,v_*)\Delta\phi\dd v_*\dd\omega,\\
 C_t(\phi,\psi)&=\tfrac12\int
       B(v-v_*,\omega)f_t(x,v)f_t(x,v_*)
                  \Delta\phi\,\Delta\psi\dd x\dd v\dd v_*\dd\omega .
\end{align*}
Its initial covariance is $\int f_0\phi\psi$.
The initial field is independent of the Gaussian noise whose covariance
is $\int_0^{s\wedge t}C_r(\phi,\psi)\dd r$.
We use the source's backward mild construction and finite-dimensional
identification~\cite[\S\S2,7--8]{U}, not an assumed path-space theorem.

\begin{theorem}[Regular-lifespan finite-family convergence]\label{thm:functional}
Under \Cref{ass:kinetic} and the imported package, for every fixed
$\phi_1,\ldots,\phi_d\in C_c^\infty(\R^6)$ the exactly centered true
hard-sphere processes converge in
$D([0,T],\R^d)$ with its Skorokhod $J_1$ topology to a continuous
Gaussian version of the inherited finite-dimensional law.
The same holds for the one fixed pasted dynamics.
\end{theorem}
\begin{proof}
Choose $\Gamma$ as required by the source's component expansion and
global cutoff estimate. Choose $L$ once large enough for its original
majorants, the squared-flux bounds, and the doubled-slab/fixed-Gaussian
reservation in \Cref{lem:calendar}. All additional bases are independent
of moment order and auxiliary accuracy. Thus all estimates below concern
one microscopic process. Accuracy cutoffs are chosen only afterward.

It suffices first to work with one fixed test, writing $X_\eps(t)$ for
its pasted fluctuation. The finite-particle paths are cadlag.
Their expectation is cadlag by dominated convergence, since the
observation sum is bounded by $\|\phi\|_\infty N$ and $\E N<\infty$.

Use a uniform mesoscopic grid of spacing $h_\eps\asymp\eps^{1/4}$ and
let $Y_\eps$ linearly interpolate the values of $X_\eps$ there.
For nodes separated by $d\geq h_\eps$, \Cref{cor:moments} gives
\[
 \E|X_\eps(t)-X_\eps(s)|^4
 \leq C[(d+r_\eps)^2+\eps^2(d+r_\eps)]\leq C'd^2 .
\]
The same bound $C|t-s|^2$ holds for the interpolation at arbitrary times:
inside a cell its fourth moment has the fractional factor
$(|t-s|/h_\eps)^4$, and across cells split into the two end pieces and
the node-to-node middle interval.
Tight initial laws are inherited. The Kolmogorov criterion therefore
makes $Y_\eps$ tight in $C([0,T])$.

Subdivide each mesoscopic cell into equal intervals of length at most
$\eps^2$ and comparable to it. There are $O(\eps^{-2})$ fine-grid points.
For each point compare $X_\eps$ with its mesoscopic left endpoint,
also including right endpoints in the comparison.
The fixed eighteenth-moment case of \Cref{cor:moments} gives
\[
 \Pp\left\{\max_{\text{fine points }t}
 |X_\eps(t)-X_\eps(\ell(t))|>\lambda\right\}
       \leq C_\lambda\eps^{9/4-2}\longrightarrow0 .
\]
In particular all neighboring mesoscopic increments vanish uniformly.
\Cref{prop:micro} with $J=2$ controls every point between fine nodes.
It follows that $\sup_{t\leq T}|X_\eps(t)-Y_\eps(t)|\to0$ in probability.
Fixed-time comparisons and the inherited finite-dimensional Gaussian
convergence identify every subsequential interpolation limit. This also
constructs a continuous version of that law. A finite product argument
gives the joint assertion for the $d$ tests.

Finally couple true and pasted flows by their common initial data.
The source cutoff event $E_\eps$ concerns disagreement anywhere in
$[0,T]$ and has probability $O(\eps^{21/20})$~\cite[\S5]{U}.
Both flows conserve $N$. Uniformly in time,
\[
 \sqrt\mu\,|\E\pi_t^{\eps,\mathrm{true}}(\phi)
                    -\E\pi_t^{\eps,p}(\phi)|
 \leq2\|\phi\|_\infty\eps^{-1}
              \E[(N/\mu)\ind_{E_\eps}]
 \leq C_\phi\eps^{1/40},
\]
using the uniform forty-second moment and
$(21/20)(41/42)-1=1/40$.
Outside $E_\eps$ the uncentered paths agree at every time.
Hence the exactly centered paths differ uniformly in probability by
something tending to zero. This transfers the finite-family weak limit.
\end{proof}

\begin{remark}
Only finite families of tests are asserted in \Cref{thm:functional}.
The separate strong-dual functional theorem, including its topology
and tightness argument, is proved in \Cref{thm:tempered}. Likewise uniform path comparison does not imply
true-flow high-moment convergence: a fluctuation of size $\eps^{-1}$ on
an event of probability $\eps^{21/20}$ has diverging variance despite
vanishing probability. None of our conclusions uses that invalid transfer.
\end{remark}

\section{The strong tempered-distribution path limit}\label{sec:tempered}

Let $\mathcal S=\mathcal S(\R^6)$ be the real Schwartz space and let
$\mathcal S'_\beta$ be its dual with the strong topology. Thus the
state-space seminorms are
\[
 p_B(u)=\sup_{\phi\in B}|u(\phi)|,
 \qquad B\subset\mathcal S\text{ bounded}.
\]
The family is directed by finite unions and positive rescaling.
We use the generalized Skorokhod $J_1$ topology on the space of
strongly cadlag paths, generated by
\begin{equation}\label{eq:tempered-jone}
 d_B(x,y)=\inf_{\lambda\in\Lambda}
 \max\left\{\sup_{t\leq T}|\lambda(t)-t|,
       \sup_{t\leq T}p_B\bigl(x(\lambda(t))-y(t)\bigr)\right\},
\end{equation}
where $\Lambda$ is the set of increasing homeomorphisms of $[0,T]$
onto itself. This is the strong-state version of
\cite[\S1, equation (1.3)]{Jakubowski}, after time rescaling.
It is not defined merely as a product of scalar path topologies.
For each $\phi\in\mathcal S$, the map
\[
 \Pi_\phi:D([0,T],\mathcal S'_\beta)\longrightarrow D([0,T],\R),
 \qquad (\Pi_\phi x)(t)=x(t)(\phi),
\]
is continuous by \eqref{eq:tempered-jone}, using the bounded set
$B=\{\phi\}$. The same observation applies to finite vectors of tests.

For either microscopic dynamics define the signed measure
\begin{equation}\label{eq:tempered-field}
 Z_t^\eps=\mu^{-1/2}\left(U_t^\eps-\E U_t^\eps\right),
 \qquad U_t^\eps=\sum_{i=1}^N\delta_{z_i(t)}.
\end{equation}
The paths use the post-collision convention of \Cref{lem:convention}.
The expectation in \eqref{eq:tempered-field} can equivalently be taken
as a Bochner expectation in the auxiliary Hilbert space used below.

\begin{theorem}[Strong tempered-distribution convergence]
\label{thm:tempered}
Under \Cref{ass:kinetic} and the imported package of
\Cref{prop:imports}, the laws of \eqref{eq:tempered-field}, for both
the true and the one fixed pasted dynamics, converge weakly in
$D([0,T],\mathcal S'_\beta)$ with topology
\eqref{eq:tempered-jone}. The common limit is a Radon probability law
supported on $C([0,T],\mathcal S'_\beta)$.
Its finite-dimensional evaluations on Schwartz tests are the centered
Gaussian field of the source, restricted from $E_0^0=C_0(\R^6)$ to
$\mathcal S$. Its covariance is
\begin{align}\label{eq:tempered-covariance}
 \E[G_t(\phi)G_s(\psi)]
 ={}&\int f_0(z)U(0,t)\phi(z)U(0,s)\psi(z)\dd z\notag\\
 &+\int_0^{s\wedge t}
 C_r\bigl(U(r,t)\phi,U(r,s)\psi\bigr)\dd r,
\end{align}
where $U(r,t)$ is the source's backward propagator and $C_r$ is the
collision covariance from \Cref{sec:functional}.
\end{theorem}

Here weak convergence tests all bounded continuous functions on the
specified path space. Neither tightness in a fixed negative-Sobolev
norm nor convergence for the uniform topology on all discontinuous
paths is asserted. On the continuous-path subspace, the topology
induced by \eqref{eq:tempered-jone} is uniform convergence in every
strong-dual seminorm. We prove all test-class, path-realization,
and law-identification steps before concluding the theorem.

\subsection{The all-Schwartz scalar estimates}

\begin{lemma}\label{lem:tempered-scalar}
For every $\phi\in\mathcal S$ the pasted scalar fields
$Z^{\eps,p}(\phi)$ are $C$-tight in $D([0,T],\R)$: they are tight
and all subsequential limits are supported on continuous paths.
At every fixed order $k\geq2$, their increment cumulants satisfy
\begin{equation}\label{eq:tempered-increments}
 \left|\kap_k\bigl(Z_t^{\eps,p}(\phi)-Z_s^{\eps,p}(\phi)\bigr)\right|
 \leq C_k q(\phi)^k\eps^{k-2}(|t-s|+r_\eps),
 \quad
 q(\phi)=\|\phi\|_\infty+\|v\cdot\nabla_x\phi\|_\infty.
\end{equation}
The choice of pasted dynamics does not depend on $\phi$ or $k$.
\end{lemma}
\begin{proof}
Both terms defining $q$ are finite, and $q$ is a continuous seminorm
on $\mathcal S$. We check explicitly why compact support was not
used in the estimates of \Cref{thm:increments,prop:micro}.

The source's marked convention permits bounded continuous observation
functions~\cite[\S6, Conventions and the assertion]{U}.
A pulse $y\mapsto y+\phi(z)$ is a measurable output pushforward.
In the component kernels, stopped variation estimates, and bounded
record factors of \Cref{prop:imports}, only the record range and
the fixed number of record-bearing labels enter. A pushforward does
not increase total variation. Normalizing a nonzero test by its
supremum norm keeps a two-pulse record in $[-2,2]$; homogeneity restores
the original scale. No integration over the support of $\phi$ occurs.
The spatial integrals instead use a birth density in each component
translation and the fixed envelopes of the other birth densities.

Put $M=\|\phi\|_\infty$ and $D=\|v\cdot\nabla_x\phi\|_\infty$.
On every designated particle line of an autonomous or retained history,
the fundamental theorem of calculus between contacts and the bounded
jump estimate give
\[
 |\phi(z(t))-\phi(z(s))|
       \leq D(t-s)+2M N_{(s,t]},
\]
where the contact count may be enlarged to include the entire diagram.
The localized tree proof uses this inequality for one increment factor
and bounds the other factors by $2M$. Its contact-count insertion costs
only a fixed polynomial in diagram size. The cyclic, initial-link,
high-complexity, and actual microscopic-error terms drop bounded marks.
Thus normalization by $q(\phi)$ proves
\eqref{eq:tempered-increments} by the same proof as
\Cref{thm:increments}; the case $q(\phi)=0$ is immediate.
The diagram bases, and hence $\Gamma,L$, are unchanged.

The count $A_I$ in \Cref{prop:micro} is independent of $\phi$.
Its conversion to an observable uses only
\[
 \osc_I\left(\eps\sum_i\phi(z_i(t))\right)
       \leq \eps DN|I|+2M\eps A_I.
\]
For the mean use the same inequality with expectations. The bound
$\E N^2\leq C\mu^2$ and the count estimates therefore give the same
microscopic modulus for every Schwartz test.
The initial CLT in~\cite[\S2, Initial ensemble]{U} includes bounded
tests, so the initial scalar laws are tight.
The moment--cumulant and two-grid interpolation proof of
\Cref{thm:functional} now applies unchanged, including its fixed
high-moment choice: the continuous interpolations are tight and their
uniform distance from $Z^{\eps,p}(\phi)$ vanishes in probability.
This proves $C$-tightness for every element of $\mathcal S$, rather
than only for a dense family of compact tests.
\end{proof}

\subsection{Identification beyond compact tests}

\begin{lemma}\label{lem:tempered-fdd}
Every finite collection of pasted Schwartz observations converges
to the corresponding source Gaussian observations.
\end{lemma}
\begin{proof}
At a fixed time $t$, take the actual unmarked normalized factorial
cumulant measures $c_{1,t}^\eps,c_{2,t}^\eps$.
The source's sharp connected bound gives
\[
 \|c_{1,t}^\eps\|_{\TV}+\mu\|c_{2,t}^\eps\|_{\TV}\leq C_t
 \qquad\text{\cite[\S6, Sharp connected bounds]{U}}.
\]
These are bounds on the unmarked measures, with constants independent
of a subsequently chosen test. For any bounded measurable $h$ the
ordinary second-cumulant identity therefore yields
\begin{equation}\label{eq:tempered-variance-cutoff}
 \E|Z_t^{\eps,p}(h)|^2
 =c_{1,t}^\eps(h^2)+\mu c_{2,t}^\eps(h\otimes h)
 \leq C_t\|h\|_\infty^2.
\end{equation}
Only finitely many fixed times are at issue, so a time-uniform
constant is unnecessary.

Choose $\chi\in C_c^\infty(\R^6)$ equal to one near zero and put
$\phi_R(z)=\chi(z/R)\phi(z)$.
For $\phi\in\mathcal S$, one has
$\phi_R\in C_c^\infty$, $\|\phi_R-\phi\|_\infty\to0$, and
$\phi_R\to\phi$ in all Schwartz seminorms.
For every finite observation vector,
\eqref{eq:tempered-variance-cutoff} and a finite sum show that the
cutoff vector approximates the original one in probability uniformly
in small $\eps$ as $R\to\infty$.
Its fixed-$R$ limit is identified by \Cref{thm:functional}.

The source's Gaussian target is a continuous linear map
$E_a^0\to L^2$ uniformly in $t$, and at $a=0$ this gives
\begin{equation}\label{eq:tempered-target-bound}
 \sup_{t\leq T}\|G_t(h)\|_{L^2}\leq C\|h\|_\infty,
 \qquad h\in C_0(\R^6)
 \quad\text{\cite[\S2, Gaussian target]{U}}.
\end{equation}
Thus the Gaussian cutoff observations converge in $L^2$.
The elementary approximation theorem for weak convergence proves
the assertion and identifies the covariance with
\eqref{eq:tempered-covariance}.
The argument used an actual unweighted microscopic covariance bound,
not a weighted remainder estimate or a variance bound for the true law.
\end{proof}

\subsection{Strong paths, exact means, and the true-flow comparison}

\begin{lemma}\label{lem:tempered-realization}
For every $\eps$, each field in \eqref{eq:tempered-field} defines a
Radon Borel probability law on $D([0,T],\mathcal S'_\beta)$.
For any fixed $s>3$, the true and pasted fields can be coupled so that
\begin{equation}\label{eq:tempered-coupling}
 \sup_{t\leq T}\|Z_t^{\eps,\mathrm{true}}-Z_t^{\eps,p}\|_{H^{-s}(\R^6)}
       \longrightarrow0\quad\hbox{in probability}.
\end{equation}
\end{lemma}
\begin{proof}
Use the ordinary inhomogeneous Sobolev space $H^{-s}(\R^6)$.
For $s>3$, Fourier transformation gives, up to the fixed normalization,
\[
 \|\delta_z\|_{H^{-s}}^2
      =c\int_{\R^6}(1+|\xi|^2)^{-s}\dd\xi=C_s^2<\infty.
\]
Dominated convergence applied to $e^{-iz\cdot\xi}$ proves that
$z\mapsto\delta_z$ is norm-continuous into this separable Hilbert space.
The finite-particle flow has finitely many collisions outside its
exceptional null set; see~\cite[Proposition 1.2 and Proposition 4.5]{DHM}.
For the pasted flow, the accepted-edge caps also give this directly
in each of the finitely many slabs. \Cref{lem:convention} fixes the
right-continuous version, and no phase jump is created by a restart
or a rejected contact. Consequently $U^\eps$ is $H^{-s}$-cadlag and
\[
             \sup_{t\leq T}\|U_t^\eps\|_{H^{-s}}\leq C_sN.
\]
Since $\E N<\infty$, its fixed-time Bochner means exist.
Dominated convergence shows that the mean is right-continuous and
has left limit $\E U_{t-}^\eps$ at each $t>0$.
Thus the exactly centered field is also $H^{-s}$-cadlag.
This argument needs no mean-continuity assertion; in fact the
fixed-time collision-null property in \Cref{lem:convention} and the
same domination show that the mean is continuous.

The source flow is measurable at each deterministic time.
Fixed-time measurability and cadlag paths make $Z^\eps$ a Borel
random variable in the Polish space $D([0,T],H^{-s})$, whose Borel
sigma-field is generated by its coordinates.
We are not invoking continuity of subtraction in $J_1$ at common jumps.
For a bounded Schwartz set $B$,
\begin{equation}\label{eq:tempered-hilbert-inclusion}
 p_B(u)\leq
 \left(\sup_{\phi\in B}\|\phi\|_{H^s}\right)\|u\|_{H^{-s}}.
\end{equation}
The bracket is finite, since $\mathcal S\to H^s$ is continuous.
Hence $H^{-s}\to\mathcal S'_\beta$ and its induced path inclusion
are continuous, the latter directly by \eqref{eq:tempered-jone}.
The pushforward of a Radon law under this continuous inclusion is
Radon: for any Borel set, approximate its inverse image from inside
by compacts and take their compact images. This proves the first claim.

For the second claim let $E_\eps$ be the source's full-time discrepancy
event, used in \Cref{thm:functional}. Outside it the uncentered paths
agree, while always
\[
 \sup_t\|U_t^{\eps,\mathrm{true}}-U_t^{\eps,p}\|_{H^{-s}}
        \leq2C_sN\ind_{E_\eps}.
\]
Subtracting the two Bochner means gives
\[
 \sup_t\|Z_t^{\eps,\mathrm{true}}-Z_t^{\eps,p}\|_{H^{-s}}
 \leq2C_s\mu^{-1/2}
       \bigl(N\ind_{E_\eps}+\E[N\ind_{E_\eps}]\bigr).
\]
The deterministic term tends to zero by the exact-centering
H\"older estimate in \Cref{thm:functional}; the random term vanishes
outside an event of probability $O(\eps^{21/20})$.
This proves \eqref{eq:tempered-coupling}.
It transfers scalar $C$-tightness and the Schwartz finite-dimensional
limits to the true flow by \eqref{eq:tempered-hilbert-inclusion}.
It transfers no true-flow high moments and asserts no Hilbert-norm
tightness of the family.
\end{proof}

\subsection{Nuclear tightness, Borel determination, and subsequences}

We now apply Mitoma's theorem to the full Schwartz test space:
on the strong dual of a nuclear Fr\'echet space, tightness of every
scalar projection is equivalent to tightness of the path laws
\cite{Mitoma}. The precise strong-topology statement also appears in
\cite[\S5.II, Theorem 5.5]{Jakubowski}.
Its hypotheses hold here because $\mathcal S$ is nuclear Fr\'echet,
\Cref{lem:tempered-realization} gives Borel path laws, and
\Cref{lem:tempered-scalar} covers every $\phi\in\mathcal S$.
Therefore the pasted laws are tight along every sequence $\eps_n\to0$.
The last paragraph of \Cref{lem:tempered-realization} permits the same
application to the true laws. We do not need a nonmetrizable Slutsky
theorem. The following details justify both subsequences and uniqueness.

Choose a countable $\mathcal S$-dense subset
$\{\phi_j:j\geq1\}\subset C_c^\infty(\R^6)$, replacing it by its
rational linear span if desired. Such a set exists by separability
and the Schwartz cutoff approximation above. These tests separate
all distributions.

We first record why countably many pairings determine the strong
state-space Borel sigma-field. Take an increasing Hilbertian scale
defining $\mathcal S$, with compact connecting maps at suitably
spaced indices; Hermite-function norms provide such a scale.
The dual is the union of the negative Hilbert spaces.
Closed integer-radius balls in these spaces, viewed in sufficiently
weaker spaces, are compact, and their images are compact in
$\mathcal S'_\beta$. They form a countable cover $\{K_n\}$ of
$\mathcal S'$.
The coordinate map $u\mapsto(u(\phi_j))_j$ is a continuous injection
into $\R^{\mathbb N}$ and is a homeomorphism on each $K_n$.
A strong Borel set intersected with $K_n$ therefore has a Borel
coordinate image. The compact image of $K_n$ is Borel in the product,
and the global injection recovers that intersection by inverse image.
Taking the countable union over $n$ proves the state-space assertion.

The same Hilbert-scale argument explains the path-space hypothesis in
\cite[Proposition 5.3]{Jakubowski}. A strongly cadlag path is bounded
on $[0,T]$ after every scalar pairing. Banach--Steinhaus gives one
dual-Hilbert norm and one radius bounding its entire range, and a
compact inclusion places the range in a compact ball of the preceding
kind. Proposition 5.3(ii) gives equality of the strong path Borel
and simple cylindrical sigma-fields. Combining it with the state-space
fact and right continuity gives
\begin{equation}\label{eq:tempered-path-borel}
 \mathcal B\bigl(D([0,T],\mathcal S'_\beta)\bigr)
 =\sigma\{x\mapsto x(q)(\phi_j):j\geq1,\ q\in Q_T\},
 \quad Q_T=([0,T]\cap\mathbb Q)\cup\{T\}.
\end{equation}
Indeed every nonterminal time is approached from above by rational
times. This is a Borel-generation statement, not just separation of
points.

For completeness, tightness supplies actual weakly convergent
subsequences even though the strong path space is not Polish.
Consider the continuous injection
\[
 J:D([0,T],\mathcal S'_\beta)\longrightarrow
 F=\prod_{j\geq1}D([0,T],\R),\qquad
 Jx=(\Pi_{\phi_j}x)_j.
\]
The product $F$ is Polish, and $J$ is a homeomorphism on every compact
subset of its domain. Given a tight sequence $P_n$, choose increasing
compact sets $L_m$ with $P_n(L_m)\geq1-2^{-m}$ for every $n$.
Ordinary Polish Prokhorov yields a subsequence for which
$J_\#P_n\Rightarrow\nu$ on $F$.
Since $J(L_m)$ is compact, portmanteau gives
$\nu(J(L_m))\geq1-2^{-m}$.
Thus $\nu$ is supported on the Borel set
$A=\bigcup_mJ(L_m)$.
The inverse of $J$ on $A$ is Borel, being continuous on each compact
piece. Its pushforward defines a probability $P$ on the strong path
space. This law is Radon: restrict first to one compact piece and use
regularity of $\nu$ there to approximate every Borel set from within
by compact sets, then let the piece increase.

If $f$ is bounded continuous on the strong path space, transport its
restriction to $L_m$ to $J(L_m)$ and extend it by Tietze to a bounded
continuous $f_m$ on $F$, without enlarging its supremum norm.
For both $P_n$ and $P$, the difference between the desired integral
and that of the extension is at most $2\|f\|_\infty2^{-m}$.
First use $J_\#P_n\Rightarrow\nu$, then let $m\to\infty$.
This proves $P_n\Rightarrow P$ in the original strong topology and
establishes a Radon subsequential limit without metrizing that topology.

\subsection{Strong continuity and the unique limit law}

\begin{proof}[Proof of \Cref{thm:tempered}]
Let $P$ be any subsequential limit obtained above, for either family.
The continuous map $\Pi_{\phi_j}$ sends it to a scalar subsequential
limit. Scalar $C$-tightness implies that $\Pi_{\phi_j}x$ is continuous
for $P$-almost every $x$.
Intersect these countably many probability-one events.
Every remaining path is already strongly cadlag, and at every $t>0$
\[
       (x(t)-x(t-))(\phi_j)=0\qquad\text{for every }j.
\]
Density of the tests implies $x(t)=x(t-)$ as distributions, for every
$t>0$ simultaneously. Equality to the strong left limit, together
with strong right continuity, proves strong continuity, including
at $T$. There is no uncountable union of null sets in this argument.
The continuous-path set is Borel: it equals
$\bigcap_j\Pi_{\phi_j}^{-1}C([0,T],\R)$ on the already cadlag space.
It is also closed, and its topology from
\eqref{eq:tempered-jone} is uniform convergence in each $p_B$
\cite[Proposition 1.6(i)]{Jakubowski}.
The latter statement also follows directly from the uniform
$p_B$-continuity of a continuous limiting path and the time changes
in \eqref{eq:tempered-jone}.

Finite-vector scalar path projections are continuous into ordinary
finite-dimensional $J_1$ spaces. Their limiting paths are continuous,
so the usual finite-dimensional evaluation theorem identifies all
their time marginals with \Cref{lem:tempered-fdd}; for the true law
use \eqref{eq:tempered-coupling} as already explained.
In particular the finite-dimensional distributions of the coordinates
in \eqref{eq:tempered-path-borel} have one prescribed Gaussian law.
They determine at most one probability measure on that Borel
sigma-field. Hence every subsequential limit of either family is the
same $P$. Tightness and this uniqueness, applied to every sequence
$\eps_n\to0$, give full convergence.
Equation \eqref{eq:tempered-covariance} follows from
\Cref{lem:tempered-fdd}. This constructs the claimed Radon, strongly
continuous realization of the source Gaussian target; it was not
assumed as an input.
\end{proof}

\begin{remark}[Scope of the distribution-valued statement]
Mitoma was applied only after scalar tightness had been proved for
every Schwartz test. Compact tests alone were used for the original
finite-dimensional identification and for a countable determining
family. The actual microscopic errors remained unweighted and were
tested only with bounded records. The source's Gaussian field is used
with its existing test-scale interpretation: we do not claim that the
collision drift maps $\mathcal S$ into $\mathcal S$, or enlarge the
source's class of solutions in its uniqueness assertion.
The Hilbert space $H^{-s}$ in \Cref{lem:tempered-realization} is a
device for individual path laws and coupling, not an asserted
uniform tightness topology.
\end{remark}

\fi
\begin{thebibliography}{9}
\bibitem{U}
OpenAI, \emph{Hard-sphere fluctuations on the regular Boltzmann lifespan},
23 September 2026, source at
\href{https://github.com/openai/math/tree/adc7f1241b42e322a6451854ab7e4b4c146bf78a/preprints/Hard-sphere-fluctuations-on-the-regular-Boltzmann-lifespan-September-23-2026/build}{openai/math},
commit \texttt{adc7f1241b42e322a6451854ab7e4b4c146bf78a}.
\bibitem{DHM}
Y. Deng, Z. Hani and X. Ma,
\emph{Long time derivation of the Boltzmann equation from hard sphere dynamics},
\href{https://arxiv.org/abs/2408.07818v3}{arXiv:2408.07818v3}, 2025.
\bibitem{BGSS}
T. Bodineau, I. Gallagher, L. Saint-Raymond and S. Simonella,
\emph{Statistical dynamics of a hard sphere gas: fluctuating Boltzmann
equation and large deviations}, Ann. of Math. \textbf{198} (2023), 1047--1201.
Proof locators refer to
\href{https://arxiv.org/html/2008.10403v2}{arXiv:2008.10403v2}.
\bibitem{BGSSeq}
T. Bodineau, I. Gallagher, L. Saint-Raymond and S. Simonella,
\emph{Long-time derivation at equilibrium of the fluctuating Boltzmann equation},
Ann. Probab. \textbf{52} (2024), 217--295;
\href{https://arxiv.org/abs/2201.04514}{arXiv:2201.04514}.
\bibitem{Mitoma}
I. Mitoma, \emph{Tightness of probabilities on $C([0,1];\mathcal S')$
and $D([0,1];\mathcal S')$}, Ann. Probab. \textbf{11} (1983), 989--999,
\href{https://doi.org/10.1214/aop/1176993447}{doi:10.1214/aop/1176993447}.
\bibitem{Jakubowski}
A. Jakubowski, \emph{On the Skorokhod topology}, Ann. Inst. H. Poincar\'e
Probab. Statist. \textbf{22} (1986), 263--285;
\href{https://www.numdam.org/item/AIHPB_1986__22_3_263_0.pdf}{original article}.
\end{thebibliography}
\end{document}
